\centerline{\Huge \bf  Первый разнобойчик}

\begin{flushright}
    \scriptsize{Хоть я и фанат челси, \\
                но галактикос $-$ сильнейший состав за всю историю футбола!}
\end{flushright}


\problem{1}
Великий маг Рубик хочет оклеить кубик несколькими бумажными параллелограммами, среди которых нет прямоугольников, без щелей и наложений. Найдите хотя бы один такой способ!

\problem{2}
При каких $a$ и $b$ уравнение $ x^3 + ax + b = 0 $ имеет три различных решения, составляющих арифметическую прогрессию?

\problem{3}
Известно, что число  $a + \dfrac{1}{a} \ -$ целое. Докажите, что число  $a^2 + \dfrac{1}{a^2} \ -$ тоже целое. 

\problem{4}
Имелось 2024 числа, ни одно из которых не равно нулю. Для каждой пары чисел записали их произведение. Докажите, что среди выписанных произведений не менее трети положительны.

\problem{5}
$\bigl[\textbf{Финал ВсОШ, 2024, 9.1}\bigr]$
Петя и Вася знают лишь натуральные числа, не превосходящие
$10^9 - 4000$. Петя считает хорошими числа, представимые в виде $abc + ab + ac + bc,$ где $a, b$ и cнатуральные числа, не меньшие $100.$ Вася считает хорошими числа, представимые в виде $xyz - x - y - z$, где $x, y$ и $z$ натуральные числа, большие $100.$ Для кого из них хороших чисел больше?

\problem{6}
$\bigl[\textbf{Финал ВсОШ, 2024, 11.6}\bigr]$
Остроугольный неравнобедренный треугольник $ABC$ вписан в окружность $\omega$ с центром в точке $O$, его высоты пересекаются в точке $H$. Через точку $O$ проведена прямая, перпендикулярная $AH$, а через точку $H$ прямая, перпендикулярная $AO$. Докажите, что точки пересечения этих прямых со сторонами $AB$ и $AC$ лежат на одной окружности, которая касается окружности $\omega$
